\begin{xltabular}{\textwidth}{| >{\bfseries}l | X |}

% --- INTESTAZIONE PRIMA PAGINA ---
\caption{Caso d'Uso: Prenotazione a Breve Termine} \label{tab:uc_prenotazione} \\
\hline
\rowcolor{gray!15} Campo & \textbf{Contenuto} \\ \hline
\endfirsthead

% --- INTESTAZIONE PAGINE SUCCESSIVE ---
\multicolumn{2}{c}{{\tablename\ \thetable{} -- Continua dalla pagina precedente}} \\
\hline
\rowcolor{gray!15} Campo & \textbf{Contenuto} \\ \hline
\endhead

% --- PIÈ DI PAGINA (se si spezza) ---
\hline
\multicolumn{2}{r}{\textit{Continua nella pagina successiva...}} \\
\endfoot

% --- PIÈ DI PAGINA FINALE ---
\endlastfoot

% --- CONTENUTO ---
Nome & Prenotazione a Breve Termine \\ \hline
Livello & Obiettivo Utente \\ \hline
Descrizione & Il Conducente riserva per 15 minuti una colonnina \texttt{ACTIVE}, per assicurarsene la disponibilità al proprio arrivo. \\ \hline
Pre-condizioni & \parbox[t]{\linewidth}{%
Il Conducente ha effettuato l'accesso.\\
Non ha già un'altra prenotazione attiva.\\
La colonnina scelta è \texttt{ACTIVE} ed è compatibile con il connettore del veicolo.
} \\ \hline
Post-condizioni & La colonnina è in stato \texttt{RESERVED}, il campo \texttt{reserved\_by\_id} punta al Conducente e \texttt{expiration\_timestamp} è impostato a 15 minuti dall'istante corrente. La colonnina resta visibile al solo Conducente che l'ha prenotata. \\ \hline
Attori & Conducente \\ \hline
Unit Test & \texttt{StationServiceTest} \\ \hline
Flusso Principale &
\begin{enumerate}[leftmargin=*, nosep]
    \item Il Conducente seleziona una colonnina \texttt{ACTIVE} dall'elenco di quelle disponibili.
    \item Conferma di volerla prenotare.
    \item Il sistema tenta l'acquisizione con un \texttt{UPDATE} condizionale, che modifica la riga solo se la colonnina risulta ancora \texttt{ACTIVE}.
    \item Il sistema porta la colonnina a \texttt{RESERVED} e ne registra il Conducente come titolare.
    \item Il sistema imposta \texttt{expiration\_timestamp} a 15 minuti dall'istante corrente.
    \item Il sistema conferma la prenotazione mostrando l'orario esatto di scadenza.
\end{enumerate} \\ \hline
Flussi Alternativi &
\textbf{3. Colonnina non più disponibile:} \newline
3.1. Se nel frattempo un altro Conducente l'ha prenotata o occupata, l'\texttt{UPDATE} non aggiorna alcuna riga e il sistema nega la prenotazione. \newline
3.2. La transazione viene annullata con un \texttt{rollback} e l'oggetto in memoria non viene modificato, così che memoria e database non divergano. \newline
3.3. Il sistema mostra il messaggio \textit{"La colonnina non è più disponibile"} e si torna al passo 1. \newline\newline
\textbf{Nota sulla concorrenza:} \newline
\textbullet\ È il database stesso a impedire la doppia prenotazione: la condizione sullo stato \texttt{ACTIVE} è parte dell'\texttt{UPDATE}, quindi fra due Conducenti simultanei ne può vincere uno solo, senza alcun blocco applicativo. \newline\newline
\textbf{Nota sulla scadenza:} \newline
\textbullet\ Se i 15 minuti trascorrono senza che venga avviata una sessione, il \texttt{GridMonitor} riporta la colonnina ad \texttt{ACTIVE} (Tab.~\ref{tab:uc_update_state}). \\ \hline

\end{xltabular}
